
\documentclass[11pt,a4paper]{article}

\usepackage[margin=1in]{geometry}
\usepackage[T1]{fontenc}
\usepackage[utf8]{inputenc}
\usepackage{lmodern}
\usepackage{microtype}
\usepackage{xcolor}
\usepackage{hyperref}
\usepackage{booktabs}
\usepackage{longtable}
\usepackage{array}
\usepackage{enumitem}
\usepackage{titlesec}
\usepackage{tabularx}
\usepackage{multirow}
\usepackage{amsmath}
\usepackage{setspace}

\hypersetup{
    colorlinks=true,
    linkcolor=blue,
    urlcolor=blue
}

\setstretch{1.08}
\setlist[itemize]{leftmargin=1.5em}
\setlist[enumerate]{leftmargin=1.8em}

\titleformat{\section}{\large\bfseries}{\thesection.}{0.6em}{}
\titleformat{\subsection}{\normalsize\bfseries}{\thesubsection.}{0.6em}{}
\titleformat{\subsubsection}{\normalsize\bfseries}{\thesubsubsection.}{0.6em}{}

\title{\textbf{Database Design Specification}\\[4pt]
for an AI Contest Management, Judging, and Practice Platform}
\author{}
\date{\today}

\begin{document}
\maketitle

\begin{abstract}
This document presents a detailed database specification for a web-based platform designed to support AI contests, virtual contests, and post-contest practice. The schema is intentionally centered around \textit{contest entries} rather than directly around teams, because a contest participant can be an individual or a team, and the same user may participate in multiple teams with different lineups across contests. The design also supports multiple tasks per contest, multiple phases per task, asynchronous evaluation jobs, official/virtual/practice modes, communication between contestants and organizers, and current leaderboard views for both task-level and contest-level ranking.
\end{abstract}

\tableofcontents
\newpage

\section{Purpose of the Database}
The goal of this database is to serve as the \textbf{system of record} for an AI competition platform. The platform must support:
\begin{itemize}
    \item official contests,
    \item virtual replays of past contests,
    \item post-contest practice submissions,
    \item individual participation and team participation,
    \item dynamic team lineups per contest,
    \item multiple tasks inside one contest,
    \item multiple phases for each task, such as public test, private test, final public, and final private,
    \item submission upload, validation, judging, and rejudging,
    \item task-level and contest-level leaderboards,
    \item announcements, clarification requests, and technical tickets.
\end{itemize}

This design favors clarity, maintainability, and practical implementation over premature generalization.

\section{Technology Choice: SQL, NoSQL, or Hybrid}
\subsection{Recommended Architecture}
The recommended architecture is:
\begin{itemize}
    \item \textbf{PostgreSQL} as the primary relational database,
    \item \textbf{JSONB} fields inside PostgreSQL for flexible configuration and structured payloads,
    \item \textbf{Redis} for caching, queues, and short-lived operational data,
    \item \textbf{MinIO or S3-compatible object storage} for uploaded files, ZIP archives, logs, and large artifacts.
\end{itemize}

\subsection{Why PostgreSQL Should Be the Core}
This platform contains many strongly related entities:
\begin{itemize}
    \item users, teams, and memberships,
    \item contests, tasks, and phases,
    \item contest entries and lineups,
    \item submissions and judging jobs,
    \item leaderboard rows tied to contest entries,
    \item clarification and ticket flows tied to contest participation.
\end{itemize}

These entities require:
\begin{itemize}
    \item foreign keys,
    \item uniqueness constraints,
    \item transactional consistency,
    \item predictable joins,
    \item strong integrity guarantees.
\end{itemize}

These are exactly the strengths of a relational database.

\subsection{Why Not Use NoSQL as the Main Database}
A NoSQL-first design would make some flexible fields easier to store, but it would complicate:
\begin{itemize}
    \item enforcing one source of truth for contest participation,
    \item guaranteeing one active lineup per contest entry,
    \item maintaining leaderboard consistency,
    \item querying official versus virtual participation,
    \item joining user/team/entry/submission/leaderboard data.
\end{itemize}

Therefore, NoSQL is not recommended as the primary database for this project.

\subsection{Why a Hybrid Approach Is Best}
A hybrid design gives the best balance:
\begin{itemize}
    \item PostgreSQL stores structured and relational domain data,
    \item JSONB stores flexible configuration and result payloads,
    \item Redis handles operational speed concerns,
    \item object storage handles large files outside the relational database.
\end{itemize}

\section{Core Design Principles}
The schema is built on the following principles:

\subsection{Contest-Centric but Entry-Driven}
The most important design decision is that the database is \textbf{not centered directly on team IDs}. Instead, it is centered on \textbf{contest entries}. A contest entry is the actual participant in a specific contest. It may represent:
\begin{itemize}
    \item one individual participant,
    \item or one team with a contest-specific lineup.
\end{itemize}

This makes the model much cleaner for:
\begin{itemize}
    \item official contests,
    \item virtual contests,
    \item practice mode,
    \item lineups that vary across contests,
    \item users belonging to multiple teams.
\end{itemize}

\subsection{Task-Phase Separation}
Each contest can contain multiple tasks, and each task can contain multiple phases. This supports workflows such as:
\begin{itemize}
    \item public test,
    \item private test,
    \item final public,
    \item final private.
\end{itemize}

Each phase can have its own judging logic, submission limit, leaderboard visibility, and availability for official/virtual/practice modes.

\subsection{One Final Score Per Submission in V1}
For the first version, each submission stores one final score:
\begin{itemize}
    \item a raw score used for ranking,
    \item a display score shown in the UI,
    \item an optional JSON breakdown for debugging or future reporting.
\end{itemize}

This keeps the system simple while remaining flexible enough for AI contests.

\subsection{Split Leaderboards by Scope}
Instead of forcing task-scoped and contest-scoped ranking into one generic structure, the database uses:
\begin{itemize}
    \item one table for task-phase leaderboard rows,
    \item one table for contest-phase leaderboard rows.
\end{itemize}

This makes the leaderboard model easier to reason about and avoids a weak abstraction based on free-form board identifiers.

\subsection{Flexible Metadata with JSONB}
Fields that can vary by contest or by judging logic are stored as JSONB:
\begin{itemize}
    \item contest rules,
    \item task submission schema,
    \item submission validation result,
    \item submission score payload,
    \item leaderboard score breakdown.
\end{itemize}

\section{High-Level Entity Relationship View}
The main relationships are:

\begin{itemize}
    \item A \textbf{user} can belong to multiple teams.
    \item A \textbf{team} can contain multiple users.
    \item A \textbf{contest} contains multiple tasks.
    \item A \textbf{task} contains multiple phases.
    \item A \textbf{contest entry} belongs to exactly one contest.
    \item A \textbf{contest entry} represents either one user or one team.
    \item A \textbf{contest entry} has one lineup through contest\_entry\_members.
    \item A \textbf{submission} belongs to one contest entry, one task, and one phase.
    \item A \textbf{submission} can have one or more stored uploaded files.
    \item A \textbf{submission} can generate one or more evaluation jobs.
    \item A \textbf{task-phase leaderboard entry} belongs to one contest entry, one task, and one real task phase.
    \item A \textbf{contest-phase leaderboard entry} belongs to one contest entry and one logical contest phase definition.
\end{itemize}

\section{Detailed Table Specification}

\subsection{Table 1: users}
\textbf{Purpose:} stores user accounts in the system.

\begin{longtable}{>{\raggedright\arraybackslash}p{0.18\textwidth} >{\raggedright\arraybackslash}p{0.17\textwidth} >{\raggedright\arraybackslash}p{0.57\textwidth}}
\toprule
\textbf{Column} & \textbf{Type} & \textbf{Meaning} \\
\midrule
id & UUID, PK & Internal primary key for the user. \\
email & VARCHAR, UNIQUE & Login email address. Must be unique. \\
password\_hash & VARCHAR & Hashed password. Never store plaintext passwords. \\
full\_name & VARCHAR & User's display name. \\
role & ENUM/VARCHAR & Role in the platform: contestant, jury, or admin. \\
student\_id & VARCHAR, NULL & Student identifier if the institution needs it. \\
avatar\_url & VARCHAR, NULL & Optional profile image URL. \\
last\_visit & TIMESTAMPTZ, NULL & Last known activity timestamp. \\
created\_at & TIMESTAMPTZ & Creation time. \\
updated\_at & TIMESTAMPTZ & Last update time. \\
\bottomrule
\end{longtable}

\textbf{Notes:}
\begin{itemize}
    \item The role belongs to the platform, not necessarily to a team.
    \item A contestant may later become jury or admin without changing their identity record.
\end{itemize}

\subsection{Table 2: teams}
\textbf{Purpose:} stores global teams, not tied directly to a specific contest.

\begin{longtable}{>{\raggedright\arraybackslash}p{0.18\textwidth} >{\raggedright\arraybackslash}p{0.17\textwidth} >{\raggedright\arraybackslash}p{0.57\textwidth}}
\toprule
\textbf{Column} & \textbf{Type} & \textbf{Meaning} \\
\midrule
id & UUID, PK & Team primary key. \\
slug & VARCHAR, UNIQUE & Short readable identifier for URLs and references. \\
name & VARCHAR & Human-readable team name. \\
owner\_id & UUID, FK users(id) & Team creator or owner. \\
created\_at & TIMESTAMPTZ & Creation time. \\
updated\_at & TIMESTAMPTZ & Last update time. \\
\bottomrule
\end{longtable}

\textbf{Notes:}
\begin{itemize}
    \item The same team can participate in multiple contests.
    \item Team composition may exist globally, but contest-specific lineups are handled separately.
\end{itemize}

\subsection{Table 3: team\_members}
\textbf{Purpose:} stores the many-to-many relation between users and teams.

\begin{longtable}{>{\raggedright\arraybackslash}p{0.18\textwidth} >{\raggedright\arraybackslash}p{0.17\textwidth} >{\raggedright\arraybackslash}p{0.57\textwidth}}
\toprule
\textbf{Column} & \textbf{Type} & \textbf{Meaning} \\
\midrule
team\_id & UUID, FK teams(id) & Team identifier. \\
user\_id & UUID, FK users(id) & User identifier. \\
role & ENUM/VARCHAR & Team role such as owner, manager, or member. \\
joined\_at & TIMESTAMPTZ & When the user joined the team. \\
\bottomrule
\end{longtable}

\textbf{Primary key:} \texttt{(team\_id, user\_id)}

\textbf{Notes:}
\begin{itemize}
    \item One user may belong to many teams.
    \item One team may contain many users.
\end{itemize}

\subsection{Table 4: contests}
\textbf{Purpose:} stores one contest event.

\begin{longtable}{>{\raggedright\arraybackslash}p{0.18\textwidth} >{\raggedright\arraybackslash}p{0.17\textwidth} >{\raggedright\arraybackslash}p{0.57\textwidth}}
\toprule
\textbf{Column} & \textbf{Type} & \textbf{Meaning} \\
\midrule
id & UUID, PK & Contest primary key. \\
slug & VARCHAR, UNIQUE & Human-readable contest code for URLs. \\
title & VARCHAR & Contest title. \\
description & TEXT & Contest description. \\
banner\_url & VARCHAR, NULL & Optional banner or poster URL. \\
status & ENUM/VARCHAR & Lifecycle state: draft, registration\_open, running, ended, archived. \\
entry\_policy & ENUM/VARCHAR & Whether the contest allows individual entries, team entries, or both. \\
registration\_start & TIMESTAMPTZ & Registration opening time. \\
registration\_end & TIMESTAMPTZ & Registration closing time. \\
start\_time & TIMESTAMPTZ & Official contest start time. \\
end\_time & TIMESTAMPTZ & Official contest end time. \\
visibility & ENUM/VARCHAR & Public or private contest visibility. \\
rules\_json & JSONB & Flexible contest rules configuration. \\
created\_by & UUID, FK users(id) & Contest creator. \\
max\_team\_size & INTEGER & Maximum allowed lineup size for team entries. \\
require\_approval & BOOLEAN & Whether entry approval is required. \\
created\_at & TIMESTAMPTZ & Creation time. \\
updated\_at & TIMESTAMPTZ & Update time. \\
\bottomrule
\end{longtable}

\subsection{Table 5: tasks}
\textbf{Purpose:} stores contest tasks. A contest may contain one or multiple tasks.

\begin{longtable}{>{\raggedright\arraybackslash}p{0.18\textwidth} >{\raggedright\arraybackslash}p{0.17\textwidth} >{\raggedright\arraybackslash}p{0.57\textwidth}}
\toprule
\textbf{Column} & \textbf{Type} & \textbf{Meaning} \\
\midrule
id & UUID, PK & Task primary key. \\
contest\_id & UUID, FK contests(id) & Parent contest. \\
slug & VARCHAR & Short task identifier, unique within a contest. \\
title & VARCHAR & Task title. \\
description & TEXT & Task summary. \\
problem\_statement\_url & VARCHAR, NULL & Link to statement or PDF. \\
submission\_schema & JSONB & Definition of expected uploaded structure. \\
score\_label & VARCHAR & UI label for the score, e.g., Score, Accuracy, MAPE. \\
higher\_is\_better & BOOLEAN & Whether larger scores rank higher. \\
sort\_order & INTEGER & Display order inside the contest. \\
created\_at & TIMESTAMPTZ & Creation time. \\
updated\_at & TIMESTAMPTZ & Update time. \\
\bottomrule
\end{longtable}

\textbf{Recommended uniqueness:} \texttt{UNIQUE(contest\_id, slug)}

\subsection{Table 6: contest\_phase\_defs}
\textbf{Purpose:} stores logical contest-wide phase definitions such as \texttt{public\_test}, \texttt{final\_public}, \texttt{private\_test}, and \texttt{final\_private}.

\begin{longtable}{>{\raggedright\arraybackslash}p{0.18\textwidth} >{\raggedright\arraybackslash}p{0.17\textwidth} >{\raggedright\arraybackslash}p{0.57\textwidth}}
\toprule
\textbf{Column} & \textbf{Type} & \textbf{Meaning} \\
\midrule
id & UUID, PK & Contest-phase definition primary key. \\
contest\_id & UUID, FK contests(id) & Parent contest. \\
key & ENUM/VARCHAR & Stable logical key such as \texttt{public\_test} or \texttt{final\_private}. \\
title & VARCHAR & Human-readable title shown in UI. \\
sort\_order & INTEGER & Display order among contest-level phase tabs. \\
\bottomrule
\end{longtable}

\textbf{Required uniqueness:} \texttt{UNIQUE(contest\_id, key)}

\textbf{Design note:}\\
This table is the source of truth for logical phase identities used consistently across tasks and contest-level leaderboards. Typical values are:
\begin{itemize}
    \item \texttt{public\_test}
    \item \texttt{final\_public}
    \item \texttt{private\_test}
    \item \texttt{final\_private}
\end{itemize}

For V1, each contest is expected to define exactly these four logical phase definitions. This completeness rule is best enforced in application or service-layer logic when creating or publishing the contest.

\subsection{Table 7: phases}
\textbf{Purpose:} stores real task-specific phases. Each phase belongs to one task and maps to one logical contest phase definition.

\begin{longtable}{>{\raggedright\arraybackslash}p{0.18\textwidth} >{\raggedright\arraybackslash}p{0.17\textwidth} >{\raggedright\arraybackslash}p{0.57\textwidth}}
\toprule
\textbf{Column} & \textbf{Type} & \textbf{Meaning} \\
\midrule
id & UUID, PK & Phase primary key. \\
task\_id & UUID, FK tasks(id) & Parent task. \\
contest\_phase\_def\_id & UUID, FK contest\_phase\_defs(id) & Logical contest-phase identity for this task phase. \\
slug & VARCHAR & Short phase identifier. \\
title & VARCHAR & Human-readable phase name. \\
description & TEXT, NULL & Optional explanation of the phase. \\
open\_time & TIMESTAMPTZ & Phase opening time. \\
close\_time & TIMESTAMPTZ & Phase closing time. \\
judge\_key & VARCHAR & Key that maps to the judging script for this phase. \\
submission\_limit & INTEGER, NULL & Maximum number of allowed submissions. \\
leaderboard\_mode & ENUM/VARCHAR & Ranking mode, for example best or latest. \\
allow\_official\_submit & BOOLEAN & Whether official entries may submit. \\
allow\_virtual\_submit & BOOLEAN & Whether virtual entries may submit. \\
allow\_practice\_submit & BOOLEAN & Whether practice entries may submit. \\
display\_scores & BOOLEAN & Whether scores are visible to contestants. \\
is\_frozen & BOOLEAN & Whether the task-phase leaderboard is frozen. \\
is\_final & BOOLEAN & Whether this phase is considered a final-model or final-result phase. \\
sort\_order & INTEGER & Display order inside the task. \\
created\_at & TIMESTAMPTZ & Creation time. \\
updated\_at & TIMESTAMPTZ & Update time. \\
\bottomrule
\end{longtable}

\textbf{Required uniqueness:}
\begin{itemize}
    \item \texttt{UNIQUE(task\_id, slug)}
    \item \texttt{UNIQUE(task\_id, contest\_phase\_def\_id)}
    \item \texttt{UNIQUE(id, task\_id)} to support composite foreign keys from child tables
\end{itemize}

\textbf{V1 rule:}\\
Each task is expected to instantiate exactly one phase for each contest-wide phase definition. In other words, if a contest defines \texttt{public\_test}, \texttt{private\_test}, \texttt{final\_public}, and \texttt{final\_private}, then every task should have exactly four corresponding rows in \texttt{phases}. The uniqueness constraint above prevents duplicates, while the completeness rule is best enforced in application or service-layer logic.

\textbf{Interpretation of judge\_key:}\\
This is not a metric. It is the identifier of the judging script for that task-specific phase. For example:
\begin{itemize}
    \item \texttt{task1\_public\_test}
    \item \texttt{task1\_final\_public}
    \item \texttt{task2\_private\_test}
\end{itemize}

\subsection{Table 8: contest\_entries}
\textbf{Purpose:} stores actual participants in a contest.

A contest entry is the unit that competes. It may be:
\begin{itemize}
    \item one individual,
    \item or one team.
\end{itemize}

It may also be in one of three modes:
\begin{itemize}
    \item official,
    \item virtual,
    \item practice.
\end{itemize}

\begin{longtable}{>{\raggedright\arraybackslash}p{0.18\textwidth} >{\raggedright\arraybackslash}p{0.17\textwidth} >{\raggedright\arraybackslash}p{0.57\textwidth}}
\toprule
\textbf{Column} & \textbf{Type} & \textbf{Meaning} \\
\midrule
id & UUID, PK & Contest entry primary key. \\
contest\_id & UUID, FK contests(id) & Parent contest. \\
entry\_type & ENUM/VARCHAR & individual or team. \\
entry\_mode & ENUM/VARCHAR & official, virtual, or practice. \\
user\_id & UUID, NULL, FK users(id) & Used when the entry is individual. \\
team\_id & UUID, NULL, FK teams(id) & Used when the entry is team-based. \\
display\_name & VARCHAR & Name shown in leaderboards. \\
status & ENUM/VARCHAR & pending, approved, active, disqualified, finished. \\
registered\_by & UUID, FK users(id) & Who created the entry. \\
approved\_by & UUID, NULL, FK users(id) & Approver if approval is required. \\
approved\_at & TIMESTAMPTZ, NULL & Approval timestamp. \\
start\_at & TIMESTAMPTZ, NULL & Used for virtual mode start time. \\
end\_at & TIMESTAMPTZ, NULL & Used for virtual mode end time. \\
created\_at & TIMESTAMPTZ & Creation time. \\
updated\_at & TIMESTAMPTZ & Update time. \\
\bottomrule
\end{longtable}

\textbf{Design note:}
This table replaces the older idea of attaching submissions directly to teams.

\textbf{Required integrity rules:}
\begin{itemize}
    \item Exactly one of \texttt{user\_id} or \texttt{team\_id} must be present.
    \item If \texttt{entry\_type = individual}, then \texttt{user\_id} must be non-null and \texttt{team\_id} must be null.
    \item If \texttt{entry\_type = team}, then \texttt{team\_id} must be non-null and \texttt{user\_id} must be null.
    \item If \texttt{entry\_mode = virtual}, then \texttt{start\_at} and \texttt{end\_at} should both be non-null and must satisfy \texttt{start\_at < end\_at}.
    \item For \texttt{official} and \texttt{practice} entries, \texttt{start\_at} and \texttt{end\_at} are usually null.
\end{itemize}

\textbf{Recommended uniqueness:}
\begin{itemize}
    \item At minimum, prevent duplicate individual entries in the same mode with \texttt{UNIQUE(contest\_id, entry\_mode, user\_id)} where \texttt{user\_id IS NOT NULL}.
    \item Prevent duplicate team entries in the same mode with \texttt{UNIQUE(contest\_id, entry\_mode, team\_id)} where \texttt{team\_id IS NOT NULL}.
\end{itemize}

\subsection{Table 9: contest\_entry\_members}
\textbf{Purpose:} stores the lineup of a contest entry.

This is especially important for team entries because a global team may have many members, but only a subset may participate in a given contest.

\begin{longtable}{>{\raggedright\arraybackslash}p{0.18\textwidth} >{\raggedright\arraybackslash}p{0.17\textwidth} >{\raggedright\arraybackslash}p{0.57\textwidth}}
\toprule
\textbf{Column} & \textbf{Type} & \textbf{Meaning} \\
\midrule
contest\_entry\_id & UUID, FK contest\_entries(id) & Parent contest entry. \\
user\_id & UUID, FK users(id) & Participating user. \\
role & ENUM/VARCHAR & leader or member for this lineup. \\
joined\_at & TIMESTAMPTZ & When this user was attached to the contest entry. \\
\bottomrule
\end{longtable}

\textbf{Primary key:} \texttt{(contest\_entry\_id, user\_id)}

\textbf{Required design clarification:}
\begin{itemize}
    \item If the parent entry is \texttt{individual}, the table should contain exactly one member row and that member should be the same as \texttt{contest\_entries.user\_id}.
    \item If the parent entry is \texttt{team}, every member listed here should be a valid member of the referenced global team, unless the platform explicitly treats contest lineups as independent snapshots.
    \item A rule such as ``one user may not join two official entries in the same contest'' is best treated as application-level validation in V1 because the required uniqueness spans parent and child tables and depends on \texttt{entry\_mode}.
\end{itemize}

\subsection{Table 10: submissions}
\textbf{Purpose:} stores every submission attempt.

\begin{longtable}{>{\raggedright\arraybackslash}p{0.18\textwidth} >{\raggedright\arraybackslash}p{0.17\textwidth} >{\raggedright\arraybackslash}p{0.57\textwidth}}
\toprule
\textbf{Column} & \textbf{Type} & \textbf{Meaning} \\
\midrule
id & UUID, PK & Submission primary key. \\
contest\_id & UUID, FK contests(id) & Explicit contest scope used to enforce cross-table consistency. \\
contest\_entry\_id & UUID, FK contest\_entries(id) & Who submitted for the contest. \\
task\_id & UUID, FK tasks(id) & Target task. \\
phase\_id & UUID, FK phases(id) & Target phase. \\
submitted\_by & UUID, FK users(id) & User who actually initiated the submission. In V1 this is assumed to be a member of the contest entry lineup. \\
status & ENUM/VARCHAR & uploaded, validating, queued, running, done, failed. \\
submitted\_at & TIMESTAMPTZ & Submission timestamp. \\
file\_count & INTEGER & Number of uploaded files linked to this submission. \\
total\_size\_bytes & BIGINT & Total uploaded size. \\
manifest\_hash & VARCHAR, NULL & Hash describing submission content structure. \\
validation\_result & JSONB, NULL & Validation details, including schema mismatches if any. \\
error\_message & TEXT, NULL & Human-readable failure reason. \\
raw\_score & NUMERIC, NULL & Internal score used for ranking. \\
display\_score & NUMERIC, NULL & Formatted score for UI display. \\
score\_payload & JSONB, NULL & Optional breakdown or structured evaluator output. \\
evaluated\_at & TIMESTAMPTZ, NULL & Evaluation completion time. \\
is\_final & BOOLEAN & Marks a final submission if the contest uses that rule. \\
rejudge\_count & INTEGER & How many times the submission has been rejudged. \\
client\_ip & VARCHAR, NULL & Request IP. \\
user\_agent & VARCHAR, NULL & Client user agent. \\
created\_at & TIMESTAMPTZ & Creation time. \\
updated\_at & TIMESTAMPTZ & Update time. \\
\bottomrule
\end{longtable}

\textbf{Important note:}\\
For V1, one submission stores one final score. This is intentional and sufficient for a Kaggle-like public display of score and rank.

\textbf{Required integrity rules:}
\begin{itemize}
    \item The parent \texttt{contest\_entry} and the target \texttt{task} must belong to the same contest.
    \item The target \texttt{phase} must belong to the selected \texttt{task}.
    \item In V1, the \texttt{submitted\_by} user must belong to the submitting entry's active lineup.
\end{itemize}

\textbf{Required relational mechanism:}
\begin{itemize}
    \item Add \texttt{UNIQUE(id, contest\_id)} on \texttt{contest\_entries}.
    \item Add \texttt{UNIQUE(id, contest\_id)} on \texttt{tasks}.
    \item Add \texttt{UNIQUE(id, task\_id)} on \texttt{phases}.
    \item Enforce \texttt{FK(contest\_entry\_id, contest\_id) -> contest\_entries(id, contest\_id)}.
    \item Enforce \texttt{FK(task\_id, contest\_id) -> tasks(id, contest\_id)}.
    \item Enforce \texttt{FK(phase\_id, task\_id) -> phases(id, task\_id)}.
    \item Enforce \texttt{FK(contest\_entry\_id, submitted\_by) -> contest\_entry\_members(contest\_entry\_id, user\_id)}.
\end{itemize}

\subsection{Table 11: submission\_files}
\textbf{Purpose:} stores uploaded files that belong to a submission.

\begin{longtable}{>{\raggedright\arraybackslash}p{0.18\textwidth} >{\raggedright\arraybackslash}p{0.17\textwidth} >{\raggedright\arraybackslash}p{0.57\textwidth}}
\toprule
\textbf{Column} & \textbf{Type} & \textbf{Meaning} \\
\midrule
id & UUID, PK & File row identifier. \\
submission\_id & UUID, FK submissions(id) & Parent submission. \\
original\_filename & VARCHAR & Original filename uploaded by the user. \\
storage\_path & VARCHAR & Object key or storage path in MinIO/S3/local storage. \\
file\_size & BIGINT & File size in bytes. \\
content\_type & VARCHAR, NULL & MIME type if known. \\
hash\_sha256 & VARCHAR, NULL & File integrity hash. \\
created\_at & TIMESTAMPTZ & Creation time. \\
\bottomrule
\end{longtable}

\textbf{Design note:}
\begin{itemize}
    \item In V1, this table is intentionally minimal and can be used simply for the main uploaded artifact.
    \item It does not need to store every extracted file inside a ZIP archive.
\end{itemize}

\subsection{Table 12: evaluation\_jobs}
\textbf{Purpose:} stores asynchronous evaluation pipeline jobs.

\begin{longtable}{>{\raggedright\arraybackslash}p{0.18\textwidth} >{\raggedright\arraybackslash}p{0.17\textwidth} >{\raggedright\arraybackslash}p{0.57\textwidth}}
\toprule
\textbf{Column} & \textbf{Type} & \textbf{Meaning} \\
\midrule
id & UUID, PK & Job identifier. \\
submission\_id & UUID, FK submissions(id) & Submission being processed. \\
job\_type & ENUM/VARCHAR & validate, judge, or rejudge. \\
status & ENUM/VARCHAR & pending, running, done, failed. \\
priority & INTEGER & Relative scheduling priority. \\
worker\_id & VARCHAR, NULL & Worker identifier that processed the job. \\
attempt\_count & INTEGER & Retry count. \\
max\_attempts & INTEGER & Maximum allowed retries. \\
input\_data & JSONB, NULL & Structured input metadata for the worker. \\
output\_data & JSONB, NULL & Structured output payload from the worker. \\
started\_at & TIMESTAMPTZ, NULL & Start time. \\
completed\_at & TIMESTAMPTZ, NULL & Completion time. \\
execution\_time\_ms & INTEGER, NULL & Runtime in milliseconds. \\
error\_log & TEXT, NULL & Internal error detail. \\
celery\_task\_id & VARCHAR, NULL & Background queue task identifier if Celery is used. \\
created\_at & TIMESTAMPTZ & Creation time. \\
\bottomrule
\end{longtable}

\textbf{Design decision:}\\
This table intentionally does not store \texttt{phase\_id} as a separate column. The evaluation phase is derived from the referenced submission, which avoids denormalization drift.

\subsection{Table 13: task\_phase\_leaderboard\_entries}
\textbf{Purpose:} stores current leaderboard rows for one task at one real task phase.

\begin{longtable}{>{\raggedright\arraybackslash}p{0.18\textwidth} >{\raggedright\arraybackslash}p{0.17\textwidth} >{\raggedright\arraybackslash}p{0.57\textwidth}}
\toprule
\textbf{Column} & \textbf{Type} & \textbf{Meaning} \\
\midrule
id & UUID, PK & Task-phase leaderboard row identifier. \\
contest\_id & UUID, FK contests(id) & Contest identifier. \\
task\_id & UUID, FK tasks(id) & Ranked task. \\
phase\_id & UUID, FK phases(id) & Ranked task-specific phase. \\
contest\_entry\_id & UUID, FK contest\_entries(id) & Ranked contest entry. \\
rank & INTEGER & Rank position. \\
score & NUMERIC & Displayed ranking score. \\
score\_breakdown & JSONB, NULL & Optional detailed subscore breakdown or cached UI payload. \\
chosen\_submission\_id & UUID, NULL, FK submissions(id) & Submission currently used for ranking. \\
entries\_count & INTEGER & Number of submissions considered for this entry in this board. \\
is\_frozen & BOOLEAN & Whether this row is part of a frozen board state. \\
is\_disqualified & BOOLEAN & Whether the row is disqualified. \\
dq\_reason & TEXT, NULL & Reason for disqualification. \\
updated\_at & TIMESTAMPTZ & Update time. \\
\bottomrule
\end{longtable}

\textbf{Required uniqueness:} \texttt{UNIQUE(phase\_id, contest\_entry\_id)}

\textbf{Required integrity rules:}
\begin{itemize}
    \item The ranked \texttt{contest\_entry} must belong to the same \texttt{contest\_id} as the leaderboard row.
    \item The selected \texttt{phase\_id} must belong to the selected \texttt{task\_id}.
    \item If \texttt{chosen\_submission\_id} is non-null, it should belong to the same \texttt{contest\_entry}, \texttt{task\_id}, and \texttt{phase\_id} as the leaderboard row.
\end{itemize}

\subsection{Table 14: contest\_phase\_leaderboard\_entries}
\textbf{Purpose:} stores current contest-level leaderboard rows for one logical contest phase definition.

\begin{longtable}{>{\raggedright\arraybackslash}p{0.18\textwidth} >{\raggedright\arraybackslash}p{0.17\textwidth} >{\raggedright\arraybackslash}p{0.57\textwidth}}
\toprule
\textbf{Column} & \textbf{Type} & \textbf{Meaning} \\
\midrule
id & UUID, PK & Contest-phase leaderboard row identifier. \\
contest\_id & UUID, FK contests(id) & Contest identifier. \\
contest\_phase\_def\_id & UUID, FK contest\_phase\_defs(id) & Ranked logical contest phase. \\
contest\_entry\_id & UUID, FK contest\_entries(id) & Ranked contest entry. \\
rank & INTEGER & Rank position. \\
score & NUMERIC & Displayed aggregate ranking score. \\
score\_breakdown & JSONB, NULL & Optional task-level breakdown or cached UI payload. \\
entries\_count & INTEGER & Number of task submissions or task-phase results considered. \\
is\_frozen & BOOLEAN & Whether this row is part of a frozen board state. \\
is\_disqualified & BOOLEAN & Whether the row is disqualified. \\
dq\_reason & TEXT, NULL & Reason for disqualification. \\
updated\_at & TIMESTAMPTZ & Update time. \\
\bottomrule
\end{longtable}

\textbf{Required uniqueness:} \texttt{UNIQUE(contest\_phase\_def\_id, contest\_entry\_id)}

\textbf{Required integrity rules:}
\begin{itemize}
    \item The ranked \texttt{contest\_entry} must belong to the same \texttt{contest\_id} as the leaderboard row.
    \item The selected \texttt{contest\_phase\_def\_id} must belong to the same \texttt{contest\_id} as the leaderboard row.
    \item Official, virtual, and practice views are derived by joining to \texttt{contest\_entries.entry\_mode}; the leaderboard table itself does not duplicate \texttt{entry\_mode}.
\end{itemize}

\textbf{Practical UI interpretation:}
\begin{itemize}
    \item Official leaderboard: join to \texttt{contest\_entries} and filter \texttt{entry\_mode = official}
    \item Virtual leaderboard: join to \texttt{contest\_entries} and filter \texttt{entry\_mode = virtual}
    \item Practice leaderboard: join to \texttt{contest\_entries} and filter \texttt{entry\_mode = practice}
\end{itemize}

\subsection{Table 15: announcements}
\textbf{Purpose:} stores organizer announcements.

\begin{longtable}{>{\raggedright\arraybackslash}p{0.18\textwidth} >{\raggedright\arraybackslash}p{0.17\textwidth} >{\raggedright\arraybackslash}p{0.57\textwidth}}
\toprule
\textbf{Column} & \textbf{Type} & \textbf{Meaning} \\
\midrule
id & UUID, PK & Announcement identifier. \\
contest\_id & UUID, FK contests(id) & Parent contest. \\
task\_id & UUID, NULL, FK tasks(id) & Optional task-specific announcement. \\
title & VARCHAR & Title of the announcement. \\
content & TEXT & Body content. \\
is\_pinned & BOOLEAN & Whether this announcement is pinned. \\
is\_public & BOOLEAN & Whether contestants can see it. \\
created\_by & UUID, FK users(id) & Author. \\
created\_at & TIMESTAMPTZ & Creation time. \\
updated\_at & TIMESTAMPTZ & Update time. \\
\bottomrule
\end{longtable}

\subsection{Table 16: clarifications}
\textbf{Purpose:} stores clarification requests and responses.

\begin{longtable}{>{\raggedright\arraybackslash}p{0.18\textwidth} >{\raggedright\arraybackslash}p{0.17\textwidth} >{\raggedright\arraybackslash}p{0.57\textwidth}}
\toprule
\textbf{Column} & \textbf{Type} & \textbf{Meaning} \\
\midrule
id & UUID, PK & Clarification identifier. \\
contest\_id & UUID, FK contests(id) & Parent contest. \\
task\_id & UUID, NULL, FK tasks(id) & Related task if applicable. \\
phase\_id & UUID, NULL, FK phases(id) & Related phase if applicable. \\
contest\_entry\_id & UUID, FK contest\_entries(id) & Entry that asked the question. \\
question & TEXT & Clarification request. \\
answer & TEXT, NULL & Organizer answer. \\
is\_public & BOOLEAN & Whether the answer is public to all contestants. \\
status & ENUM/VARCHAR & pending, answered, or closed. \\
asked\_by & UUID, FK users(id) & User who asked. \\
answered\_by & UUID, NULL, FK users(id) & User who answered. \\
answered\_at & TIMESTAMPTZ, NULL & Answer timestamp. \\
created\_at & TIMESTAMPTZ & Creation time. \\
updated\_at & TIMESTAMPTZ & Update time. \\
\bottomrule
\end{longtable}

\subsection{Table 17: tickets}
\textbf{Purpose:} stores technical support issues and operational complaints.

\begin{longtable}{>{\raggedright\arraybackslash}p{0.18\textwidth} >{\raggedright\arraybackslash}p{0.17\textwidth} >{\raggedright\arraybackslash}p{0.57\textwidth}}
\toprule
\textbf{Column} & \textbf{Type} & \textbf{Meaning} \\
\midrule
id & UUID, PK & Ticket identifier. \\
submission\_id & UUID, NULL, FK submissions(id) & Related submission if applicable. \\
contest\_entry\_id & UUID, FK contest\_entries(id) & Entry that reported the issue. \\
category & ENUM/VARCHAR & upload, judge, score, or system. \\
subject & VARCHAR & Ticket subject. \\
description & TEXT & Full issue description. \\
status & ENUM/VARCHAR & open, in\_progress, resolved, rejected. \\
priority & ENUM/VARCHAR & e.g. low, normal, high, urgent. \\
assigned\_to & UUID, NULL, FK users(id) & Staff member handling the ticket. \\
created\_by & UUID, FK users(id) & Reporter. \\
created\_at & TIMESTAMPTZ & Creation time. \\
resolved\_at & TIMESTAMPTZ, NULL & Resolution time. \\
updated\_at & TIMESTAMPTZ & Update time. \\
\bottomrule
\end{longtable}

\section{Important Design Decisions Explained}

\subsection{Why We Use contest\_entries}
Without \texttt{contest\_entries}, the system would need to attach submissions and leaderboard records directly to either users or teams. That approach becomes messy when:
\begin{itemize}
    \item a contest allows both individuals and teams,
    \item a team has a different lineup in each contest,
    \item a user belongs to multiple teams,
    \item virtual and practice participation must be stored separately from official participation.
\end{itemize}

The \texttt{contest\_entries} table solves all of these in one abstraction.

\subsection{Why We Store Submission Scores Directly in submissions}
For V1, the UI only needs one final score per submission, similar to many competition platforms where contestants mostly care about rank and score. Therefore:
\begin{itemize}
    \item we do not need a separate \texttt{scores} table yet,
    \item we do not need a separate \texttt{metrics} table yet,
    \item we only store \texttt{raw\_score}, \texttt{display\_score}, and \texttt{score\_payload}.
\end{itemize}

This keeps the schema simple without preventing future extension.

\subsection{Why We Split Leaderboards into Two Tables}
The final schema intentionally uses two leaderboard tables:
\begin{itemize}
    \item \texttt{task\_phase\_leaderboard\_entries} for one task at one real task phase,
    \item \texttt{contest\_phase\_leaderboard\_entries} for one contest-wide logical phase.
\end{itemize}

This is clearer than one generalized table because:
\begin{itemize}
    \item there is no \texttt{board\_type} column to interpret,
    \item there is no free-form \texttt{phase\_key},
    \item task-scoped ranking and contest-scoped ranking have different natural keys.
\end{itemize}

\subsection{Why submission\_files Is Kept Even in a Simple Version}
Although the first version may only upload one main ZIP file or one prediction file, keeping \texttt{submission\_files} is still useful because:
\begin{itemize}
    \item it preserves file metadata cleanly,
    \item it allows future multi-file submissions,
    \item it separates file storage concerns from the submission record itself.
\end{itemize}

\section{Suggested Integrity Constraints}
The following constraints are strongly recommended:

\begin{itemize}
    \item \texttt{users.email} must be unique.
    \item \texttt{teams.slug} must be unique.
    \item \texttt{contests.slug} must be unique.
    \item \texttt{CHECK(registration\_start <= registration\_end)} on contests when both values are present.
    \item \texttt{CHECK(start\_time < end\_time)} on contests.
    \item \texttt{CHECK(max\_team\_size > 0)} on contests.
    \item \texttt{UNIQUE(contest\_id, key)} on \texttt{contest\_phase\_defs}.
    \item \texttt{UNIQUE(contest\_id, slug)} on tasks.
    \item \texttt{UNIQUE(id, contest\_id)} on tasks.
    \item \texttt{UNIQUE(task\_id, slug)} on phases.
    \item \texttt{UNIQUE(task\_id, contest\_phase\_def\_id)} on phases.
    \item \texttt{UNIQUE(id, task\_id)} on phases.
    \item \texttt{CHECK(open\_time < close\_time)} on phases.
    \item \texttt{CHECK(submission\_limit IS NULL OR submission\_limit >= 0)} on phases.
    \item \texttt{PRIMARY KEY(team\_id, user\_id)} on team\_members.
    \item \texttt{UNIQUE(id, contest\_id)} on contest\_entries.
    \item A \texttt{CHECK} on \texttt{contest\_entries} to enforce exactly one participant source: either \texttt{user\_id} or \texttt{team\_id}.
    \item \texttt{PRIMARY KEY(contest\_entry\_id, user\_id)} on contest\_entry\_members.
    \item A relational consistency rule ensuring each submission references a \texttt{contest\_entry}, \texttt{task}, and \texttt{phase} from the same contest/task chain.
    \item \texttt{FK(contest\_entry\_id, contest\_id) -> contest\_entries(id, contest\_id)} on submissions.
    \item \texttt{FK(task\_id, contest\_id) -> tasks(id, contest\_id)} on submissions.
    \item \texttt{FK(phase\_id, task\_id) -> phases(id, task\_id)} on submissions.
    \item \texttt{FK(contest\_entry\_id, submitted\_by) -> contest\_entry\_members(contest\_entry\_id, user\_id)} on submissions.
    \item \texttt{CHECK(file\_count >= 0)}, \texttt{CHECK(total\_size\_bytes >= 0)}, and \texttt{CHECK(rejudge\_count >= 0)} on submissions.
    \item \texttt{CHECK(attempt\_count >= 0)}, \texttt{CHECK(max\_attempts >= 0)}, and \texttt{CHECK(priority >= 0)} on evaluation\_jobs.
    \item \texttt{UNIQUE(phase\_id, contest\_entry\_id)} on \texttt{task\_phase\_leaderboard\_entries}.
    \item \texttt{UNIQUE(contest\_phase\_def\_id, contest\_entry\_id)} on \texttt{contest\_phase\_leaderboard\_entries}.
\end{itemize}

\subsection{What Must Be Enforced by the Database}
The following rules are strong candidates for database-level enforcement rather than only application-level validation:
\begin{itemize}
    \item contest entry participant exclusivity,
    \item contest consistency between entries, tasks, phases, and submissions,
    \item parent-child phase/task consistency,
    \item non-negative counters and valid time ranges,
    \item uniqueness of leaderboard rows for each task-phase scope and contest-phase scope.
\end{itemize}

\subsection{What May Remain in Application Logic in V1}
The following rules may be validated in application code first if they are too awkward to express in pure SQL, though they should still be documented explicitly:
\begin{itemize}
    \item whether each contest defines exactly the required four logical phase definitions,
    \item whether each task instantiates exactly one phase for each required contest phase definition,
    \item whether a team-entry lineup user must already belong to the global team,
    \item whether one user may join multiple contest entries in the same contest under special rules,
    \item whether practice or virtual entries use different temporal submission windows beyond the base phase windows,
    \item whether a chosen submission stored in a leaderboard row is checked purely in SQL or verified before upsert in application logic.
\end{itemize}

\section{Suggested Indexes}
The following indexes are recommended for performance:

\begin{itemize}
    \item \texttt{users(email)}
    \item \texttt{teams(slug)}
    \item \texttt{contests(slug, status, start\_time, end\_time)}
    \item \texttt{contest\_phase\_defs(contest\_id, key, sort\_order)}
    \item \texttt{tasks(contest\_id, sort\_order)}
    \item \texttt{phases(task\_id, open\_time, close\_time)}
    \item \texttt{contest\_entries(contest\_id, entry\_mode, status)}
    \item \texttt{contest\_entry\_members(contest\_entry\_id, user\_id)}
    \item \texttt{submissions(contest\_id, contest\_entry\_id, task\_id, phase\_id, submitted\_at)}
    \item \texttt{submissions(status)}
    \item \texttt{evaluation\_jobs(status, created\_at)}
    \item \texttt{task\_phase\_leaderboard\_entries(phase\_id, rank)}
    \item \texttt{contest\_phase\_leaderboard\_entries(contest\_phase\_def\_id, rank)}
    \item \texttt{clarifications(contest\_id, status)}
    \item \texttt{tickets(contest\_entry\_id, status)}
\end{itemize}

\section{Operational Workflow Mapped to the Database}

\subsection{Team and Contest Registration}
\begin{enumerate}
    \item Users exist in \texttt{users}.
    \item Global teams are defined in \texttt{teams}.
    \item Team membership is defined in \texttt{team\_members}.
    \item A contestant creates a \texttt{contest\_entry}.
    \item If the entry is team-based, its lineup is stored in \texttt{contest\_entry\_members}.
\end{enumerate}

\subsection{Submission and Judging}
\begin{enumerate}
    \item A contestant uploads a submission, creating a row in \texttt{submissions}.
    \item Uploaded artifacts are recorded in \texttt{submission\_files}.
    \item Validation and judging jobs are added to \texttt{evaluation\_jobs}.
    \item After judging, score fields in \texttt{submissions} are updated.
    \item Task-level leaderboard rows in \texttt{task\_phase\_leaderboard\_entries} are recomputed or updated.
    \item Contest-level leaderboard rows in \texttt{contest\_phase\_leaderboard\_entries} are recomputed or updated.
\end{enumerate}

\subsection{Communication During Contest}
\begin{itemize}
    \item General notices use \texttt{announcements}.
    \item Rule and statement questions use \texttt{clarifications}.
    \item Technical issues and complaints use \texttt{tickets}.
\end{itemize}

\section{Future Extensions}
This schema is intentionally suitable for V1, but it leaves room for extension:

\subsection{Possible V2 Additions}
\begin{itemize}
    \item \texttt{audit\_logs} for full immutable change tracking,
    \item \texttt{user\_stats\_cache} for profile summaries,
    \item \texttt{leaderboard\_snapshots} for historical ranking snapshots,
    \item a separate \texttt{scores} table if multi-metric public reporting becomes necessary,
    \item a separate \texttt{evaluators} table if the platform later supports built-in and custom judge modules as first-class data.
\end{itemize}

\subsection{Why They Are Not in V1}
They are not strictly necessary for the first functional release. The current schema already supports:
\begin{itemize}
    \item official contests,
    \item virtual contests,
    \item practice submissions,
    \item team lineups,
    \item asynchronous judging,
    \item current leaderboards,
    \item operational communication.
\end{itemize}

\section{Final Recommendation}
The recommended production-ready database foundation for the initial version is:
\begin{itemize}
    \item \textbf{PostgreSQL} as the primary relational database,
    \item \textbf{JSONB} for flexible configuration and structured payloads,
    \item \textbf{Redis} for queueing and caching,
    \item \textbf{MinIO/S3} for submission artifacts.
\end{itemize}

The schema should contain the following seventeen core tables:
\begin{enumerate}
    \item users
    \item teams
    \item team\_members
    \item contests
    \item contest\_phase\_defs
    \item tasks
    \item phases
    \item contest\_entries
    \item contest\_entry\_members
    \item submissions
    \item submission\_files
    \item evaluation\_jobs
    \item task\_phase\_leaderboard\_entries
    \item contest\_phase\_leaderboard\_entries
    \item announcements
    \item clarifications
    \item tickets
\end{enumerate}

This design is detailed enough for implementation, clean enough for maintenance, and flexible enough for future evolution into a stronger AI contest platform.

\end{document}
